% !TeX root = ../guide
\chapter{GIMP: Preparing Illustrative and Equipment Photographs}\label{ch:gimp}
	\section{Introduction}\label{sec:GIMPIntroduction}
		While ImageJ is the strict standard for analyzing scientific data and micrographs, a thesis also requires standard photography. 
		You will often need to include pictures of your experimental setup, custom-built test rigs, or field equipment. 

		\textbf{GIMP} (GNU Image Manipulation Program) is a free, open-source raster graphics editor that serves as a highly capable alternative to Adobe Photoshop. 
		This chapter will cover how to use GIMP to clean up distracting backgrounds, correct lighting, and ensure your illustrative photographs are formatted at the correct resolution for high-quality printing in your \LaTeX\ document.

		\warning{%
			You must never use GIMP's advanced cloning, healing, or manipulation tools on scientific data (such as altering a material fracture surface or removing an \enquote{outlier} particle from a microscope image). 
			GIMP should be strictly reserved for illustrative photographs where the physical appearance of the equipment, rather than extractable data, is the focus.}

	\section{Key Features of GIMP}
		GIMP provides a comprehensive suite of tools for photo retouching and image composition.

		\subsection{Layers and Masks}
			Unlike basic photo viewers, GIMP utilizes a layer-based workflow. 
			This allows you to stack edits on top of your original image without destroying it. 
			More importantly, it allows you to use \enquote{Layer Masks} to hide distracting backgrounds. 
			If you take a photo of a piece of equipment in a cluttered workshop, you can use masks to isolate the equipment and place it on a clean, professional white background for your thesis.

		\subsection{Perspective and Lens Correction}
			When taking photos of control panels or large equipment in tight laboratory spaces, the image is often skewed by lens distortion or captured at an angle. 
			GIMP includes powerful perspective transformation tools that allow you to \enquote{square up} your photos, making them look as though they were taken perfectly straight-on.

		\subsection{Color and Lighting Adjustments}
			Lab environments often have terrible fluorescent lighting that washes out the contrast in your photos. 
			GIMP's \enquote{Levels} and \enquote{Curves} tools allow you to correct the white balance and pull detail out of dark shadows, ensuring your equipment is clearly visible when printed on paper.

	\section{Getting Started with GIMP}
		GIMP is a fully-featured program with a layout that will be familiar to anyone who has used legacy versions of Photoshop.

		\subsection{Installation}
			GIMP is completely free and available for Windows, macOS, and Linux.
			\begin{enumerate}
				\item Visit the official website: \url{https://www.gimp.org/}.
				\item Click the `Download' button and select the installer for your operating system.
				\item Follow the standard installation prompts.
			\end{enumerate}

		\subsection{The Main Interface}
			When you launch GIMP, the interface is divided into three main areas, as shown in \Cref{fig:GimpMain}.
			\begin{figure}[htbp]
				\centering
				\includegraphics[width=0.8\linewidth]{example-image}
				\caption{GIMP Main Interface showing the Toolbox, Image Canvas, and Layers Dialog.}
				\label{fig:GimpMain}
			\end{figure}
			
			\begin{itemize}
				\item \textbf{Toolbox (Left):} Contains all your selection, drawing, and transformation tools.
				\item \textbf{Image Canvas (Center):} Your primary workspace.
				\item \textbf{Dockable Dialogs (Right):} This panel houses your Layers, Channels, and Undo History.
			\end{itemize}

	\section{GIMP in a \LaTeX\ Workflow}
		The most important function GIMP serves in your thesis workflow is managing image resolution and file formats before they ever touch your \path{.tex} files.

		\subsection{Managing Print Resolution (DPI)}
			A common issue when writing a thesis is that photos taken on a smartphone look great on a screen but print out pixelated and blurry. 
			This is an issue of Dots Per Inch (DPI). 
			
			Before exporting an image for \LaTeX, you must ensure it has a high enough resolution for physical printing:
			\begin{enumerate}
				\item Open your photo in GIMP.
				\item Go to `Image' \(\rightarrow\) `Print Size'.
				\item Ensure the X and Y resolution are set to a minimum of \textbf{300 pixels/in}. 
				\item The dialog will show you the physical size (in inches or millimeters) the image will be when printed at that high resolution.
			\end{enumerate}

		\subsection{Cropping for Document Flow}
			\LaTeX\ handles text wrapping and figure placement based strictly on the boundary of the image file. 
			If your photograph has a massive amount of empty floor or ceiling space, \LaTeX\ will include that dead space, pushing your text further down the page. 
			Always use GIMP's Crop tool (`Shift + C') to trim your photographs as tightly as possible to the subject matter before exporting.

		\subsection{Exporting Images}
			Unlike standard programs where you click \enquote{Save,} GIMP uses `File' \(\rightarrow\) `Save' to create a \path{.xcf} project file (which \LaTeX\ cannot read). 
			To generate an image for your thesis, you must use `File' \(\rightarrow\) `Export As'.

			\begin{itemize}
				\item \textbf{PNG (\texttt{.png}):} Use this format if you have added any text annotations, arrows, or if you have removed the background to make it transparent. 
					PNG prevents the blurry artifacts that JPEG compression causes around sharp text.
				\item \textbf{JPEG (\texttt{.jpg}):} Use this format for standard, illustrative photographs of equipment or field work. 
					It provides good compression for complex colors.
			\end{itemize}